% Experimental evaluation. Packages: booktabs, multirow, amsmath. Numbers: docs/RESULTS.md.

\section{Experimental Evaluation}
\label{sec:experiments}

\subsection{Setup}
\label{subsec:setup}

\paragraph{Datasets and protocol.}
We use three CP4IM datasets: chess (minimum frequency $0.5$), mushroom ($0.1$) and soybean ($0.044$, the
threshold of the original experiments). Every run performs $T = 20$ iterations with queries of $k = 10$
patterns sampled by Flexics, item features, and the SCD learner with the defaults of the published
implementation (initial weights $1$, linear folding, $\eta = 0.17$, tilt $10$).

\paragraph{Simulated users.}
Besides the surprisingness user of the original evaluation, we use \emph{combination users}
(Section~\ref{…}). For each dataset, two users are drawn at random (seeded) among pairs of disjoint item pairs
that are each contained in $5$--$25\%$ of the test patterns; the draw depends on the data only. Each
configuration is run with $5$ seeds per user, giving $10$ runs per dataset and $30$ in total.

\paragraph{Evaluation.}
A learned model is evaluated on $1000$ random frequent itemsets of length $2$--$10$, grown one item at a
time, from which the patterns shown during the run are removed. We report pairwise accuracy (\emph{acc}),
the accuracy on the pairs whose number of liked combinations differs (\emph{acc-combo}, which ignores the
surprisingness tie-break), and tie-aware precision among the top $10\%$ (\emph{P@10\%}). Significance is
assessed with a paired sign-flip permutation test over (dataset, user, seed). As references we report the
score $|X|$ (pattern length) and two offline models fitted on all pairs of $2000$ other random patterns: one
over items only, and one over items plus the true liked pairs.

\subsection{Fidelity to the Published Implementation}
\label{subsec:fidelity}

Every correction of Section~\ref{…} can be undone individually. With all of them undone, our implementation
reproduces a run of the published code exactly (same queries, sub-patterns and learned weights).
Table~\ref{tab:fidelity} measures each correction on the surprisingness user.

\begin{table}[t]
\centering
\caption{Published implementation vs.\ corrections (surprisingness user, chess and soybean, 10 runs).
$p$: paired test against the published implementation.}
\label{tab:fidelity}
\small
\begin{tabular}{lccc}
\toprule
Configuration & acc (\%) & $p$ & P@10\% (\%) \\
\midrule
Score $= |X|$ (chess / soybean)            & 64.4 / 66.4 & --    & 26.0 / 17.0 \\
LetSIP                                      & 71.8 & 0.10  & 32.9 \\
\textsc{DiSPaLe}, published                 & 71.1 & --    & 30.7 \\
\quad + stale inner products fixed          & 71.7 & 0.002 & 33.2 \\
\quad + transferred-weight index fixed      & 71.9 & 0.002 & 33.5 \\
\quad + sampling with the learned model     & 71.3 & 0.27  & 31.6 \\
\quad + integer ICV means fixed             & 71.1 & 0.66  & 30.5 \\
\quad + in-place covers fixed               & 69.4 & 0.004 & 27.6 \\
\textsc{DiSPaLe}, all corrections           & 72.0 & 0.020 & 32.5 \\
\quad initial weights $0$ (linear folding)  & 66.9 & 0.004 & 32.5 \\
\bottomrule
\end{tabular}
\end{table}

The corrections together improve accuracy by about one point ($p = 0.02$), mainly through the stale inner
products and the wrong index of the transferred weight. Starting from zero weights, as stated in the paper,
costs five points, so we keep the published initialisation. Learning adds only $5$--$7$ points over ranking
by pattern length on this user.

\subsection{Multiple Sub-patterns with Folding}
\label{subsec:main_results}

\begin{table}[t]
\centering
\caption{Combination users, all datasets (30 runs per configuration). Fold: the published scheme; keep:
permanent sub-pattern features. $p$: paired test of acc-combo against \textsc{DiSPaLe} with
$k_{disc} = 1$ (fold).}
\label{tab:main}
\small
\begin{tabular}{llcccc}
\toprule
& Configuration & acc & acc-combo & P@10\% & $p$ \\
\midrule
& LetSIP                                   & 77.0 & 91.9 & 78.1 & 0.75 \\
\midrule
\multirow{6}{*}{fold}
& $k_{disc} = 1$                           & 77.0 & 92.0 & 78.4 & --   \\
& $k_{disc} = 3$, Top                      & 77.0 & 92.1 & 78.2 & 0.78 \\
& $k_{disc} = 3$, MMR                      & 77.0 & 92.0 & 78.3 & 0.94 \\
& $k_{disc} = 3$, coverage                 & 77.0 & 92.0 & 78.4 & 0.93 \\
& $k_{disc} = 3$, \textsc{PartitionGain}   & 77.0 & 92.0 & 78.4 & 0.99 \\
& $k_{disc} = 3$, pair coverage            & 77.0 & 92.0 & 78.4 & 0.95 \\
\midrule
\multirow{4}{*}{keep}
& $k_{disc} = 1$                           & 78.5 & 94.6 & 86.4 & $<0.001$ \\
& $k_{disc} = 3$, Top                      & 78.8 & 95.3 & 88.3 & $<0.001$ \\
& $k_{disc} = 3$, \textsc{PartitionGain}   & 78.5 & 94.9 & 86.5 & $<0.001$ \\
& $k_{disc} = 3$, \textsc{PartitionGain}, 10 sweeps & 78.7 & 96.0 & 90.0 & $<0.001$ \\
\bottomrule
\end{tabular}
\end{table}

With the fold-and-discard scheme, no selection rule changes the learned model (Table~\ref{tab:main}, fold):
all configurations reach the same accuracy as $k_{disc} = 1$ and LetSIP. The learned sub-pattern weights
remain between $0.01$ and $0.06$, against about $0.9$ for the item weights, and the model before folding
coincides with the model after folding. Running ten SCD sweeps per iteration raises the sub-pattern
weights to $0.16$--$0.44$, but the folded model still does not improve, and the additional sweeps reduce
accuracy by about $1.5$ points for every method by overfitting the pairs of each query. Since folding
returns an item-only model, it cannot represent a preference for items \emph{together}: on all datasets,
the folded models stay below the offline items-only reference on acc-combo.

\subsection{Keeping the Sub-patterns}
\label{subsec:keep_results}

When every selected sub-pattern is kept as a permanent feature, the model improves significantly on every
dataset (Table~\ref{tab:main}, keep, and Table~\ref{tab:keep_datasets}): $+1.5$ points of accuracy, $+2.6$
of acc-combo and $+8$ of P@10\% for $k_{disc} = 1$, and up to $+4.0$ and $+11.6$ with $k_{disc} = 3$ and
ten sweeps. The kept models exceed the offline items-only reference on acc-combo and approach the reference
that is given the true liked pairs. Selecting several sub-patterns per iteration adds a smaller gain on top
of keeping them: against $k_{disc} = 1$ with keep, Top ($p = 0.024$ on acc-combo) and
\textsc{PartitionGain} with ten sweeps ($p = 0.019$ on acc-combo, $p = 0.005$ on P@10\%) improve, whereas
\textsc{PartitionGain} with one sweep does not ($p = 0.25$). After Holm's correction over these six tests,
only the P@10\% improvement of \textsc{PartitionGain} with ten sweeps remains significant.

\begin{table}[t]
\centering
\caption{acc-combo / P@10\% per dataset (10 runs each).}
\label{tab:keep_datasets}
\small
\begin{tabular}{lccc}
\toprule
Configuration & chess & mushroom & soybean \\
\midrule
Score $= |X|$                                         & 84.3 / 67.5 & 79.5 / 50.0 & 80.4 / 50.5 \\
LetSIP                                                & 94.0 / 93.4 & 90.7 / 68.5 & 91.0 / 72.5 \\
Fold, $k_{disc} = 1$                                  & 94.5 / 93.6 & 90.5 / 67.7 & 91.0 / 73.8 \\
Fold, $k_{disc} = 3$, \textsc{PartitionGain}          & 94.2 / 93.3 & 90.6 / 68.8 & 91.2 / 73.1 \\
Keep, $k_{disc} = 1$                                  & 97.2 / 98.8 & 92.8 / 74.2 & 93.9 / 86.2 \\
Keep, $k_{disc} = 3$, Top                             & 97.8 / 99.3 & 93.8 / 79.3 & 94.3 / 86.4 \\
Keep, $k_{disc} = 3$, \textsc{PartitionGain}, 10 sweeps & 98.4 / 99.4 & 93.6 / 78.9 & 96.1 / 91.8 \\
\midrule
Offline, items only                                   & 94.9 / 99.5 & 91.0 / 69.0 & 93.0 / 78.0 \\
Offline, items + true liked pairs                     & 99.3 / 100  & 97.1 / 87.0 & 98.8 / 96.0 \\
\bottomrule
\end{tabular}
\end{table}

\subsection{Cost and Diversity}
\label{subsec:cost}

Measured with runs executed one at a time (chess), an iteration takes about $2.6$\,s, almost entirely spent
sampling the query; the sub-pattern search takes about $35$\,ms over about $70$ distinct candidate covers,
every selection rule less than $3$\,ms, and learning and folding about $11$\,ms. Selecting several
sub-patterns therefore adds no measurable cost at $k = 10$. Over the last ten iterations, all
configurations show the same number of distinct items and the same similarity between shown patterns.

\subsection{Discussion}
\label{subsec:discussion}

Selecting several complementary sub-patterns does not improve preference learning as long as the
sub-pattern features are folded into item weights after every iteration: the learner barely uses them,
and the folded model cannot represent combinations of items. The bottleneck is therefore the transfer
scheme, not the choice of sub-patterns. Once the sub-patterns are kept as features, they improve the model
markedly, and several sub-patterns per iteration add a further, smaller gain. Under this scheme, the plain
Top rule over distinct covers performs at least as well as partition gain, which suggests that, when
features persist across iterations, covering more candidate combinations matters more than
complementarity within a single query. These conclusions are limited to $k = 10$, $20$ iterations, item
features and synthetic combination users, whose preference has exactly the structure sub-pattern features
can express.
